
\section{Recursos humanos}
\label{sc:RH}

\subsection{Equipo de Trabajo}

El equipo de trabajo necesario para este proyecto de investigación está compuesto por los siguientes roles y personal:

\begin{itemize}
    \item \textbf{Profesor Guía:} Responsable de guiar al investigador principal en sus dudas generales del acerca del tópico del proyecto de investigación.
    \begin{itemize}
        \item \textit{Nombre del profesor guía:} Dr. Rodrigo Olivares.
    \end{itemize}

    \item \textbf{Investigador Principal:} Responsable de idear la propuesta de solución al problema planteado de investigación aplicada, toma de decisiones estratégicas, planificación y desarrollo de todas las fases de la investigación.
    \begin{itemize}
        \item \textit{Nombre del investigador principal:} Sebastián M. Guzmán.
    \end{itemize}

\end{itemize}






\section{Recursos materiales}
\label{sc:RM}

\subsection{Hardware (HW)}
El hardware necesario para este trabajo de investigación incluye:
\begin{itemize}
    \item \textbf{Computador de Alto Rendimiento (HPC):} Equipadas con procesadores de última generación (mínimo Intel Core i7 o AMD Ryzen 7), 16GB de RAM y discos duros SSD de al menos 512GB.
    \item \textbf{Servidor:} Servidor dedicado al almacenamiento y procesamiento de datos con capacidades superiores al HPC (en caso de ser necesario).
\end{itemize}

\subsection{Software (SW)}
El software requerido incluye:
\begin{itemize}
    \item \textbf{Sistemas Operativos:} Windows 11 Home, Linux (Ubuntu 24.04 LTS) o macOS 14.
    \item \textbf{Software de desarrollo de código:} Python (con bibliotecas como NumPy, Pandas, SciPy, Sklearn o respositorios Github).
    \item \textbf{Herramientas de Desarrollo:} Entornos de desarrollo integrados (IDE) como Visual Studio Code, y Jupyter Notebook (o Colab de Google).
    \item \textbf{Software de Gestión de Proyectos:} Microsoft Project o Project Libre.
\end{itemize}

\subsection{Bibliotecas}
Para la consulta de literatura y referencias, se requiere acceso a:
\begin{itemize}
    \item \textbf{Biblioteca Digital:} Acceso las bases de datos académicas como ScienceDirect (Elsevier), SpringerLink, y Web of Science (WoS) para obtener libros electrónicos, artículos científicos, y publicaciones digitales relevantes al área de investigación.
\end{itemize}
